\chapter{绝对定量、电荷积分与证据门槛}
\section{灵敏度的定义}
在式~\eqref{eq:demsfull} 中，$\kappa_j$ 把电极源通量转换为参考质量通道的等效电流。它包括采样和收集效率，单位为 C/mol。若某文献或仪器把灵敏度定义为电流对分压的比值，必须另有分压到源通量的模型，不能直接把该系数填入本软件的 $\kappa$ 字段。

对已知标定速率 $r_j^{\rm cal}$、完整响应窗和单位质量传输核，有
\begin{equation}
\kappa_j=\frac{\int x_j^{\rm cal}(t)\,\mathrm dt}
{\int r_j^{\rm cal}(t)\,\mathrm dt}.
\end{equation}
参考通道应满足 $S_{m_jj}=1$；否则幅值在 $S$ 与 $\kappa$ 之间重复计入。核质量也必须一致，否则产物总量随核的任意缩放改变。

\section{由源信号到物质量与法拉第效率}
在标定条件匹配时，
\begin{equation}
 r_j(t)=\frac{x_j(t)}{\kappa_j},\qquad
 N_j=\frac1{\kappa_j}\int_{t_a}^{t_b}x_j(t)\,\mathrm dt.
\end{equation}
若生成每摩尔产物对应 $z_j$ 摩尔电子，则产物电荷为
\begin{equation}
 Q_j=z_jFN_j.
\end{equation}
对还原过程定义阴极电荷的正值大小
\begin{equation}
 Q^- =\int_{t_a}^{t_b}\max[-I(t),0] \mathrm dt,
\qquad \mathrm{FE}_j=Q_j/Q^-.
\end{equation}
氧化过程使用 $Q^+=\int\max[I(t),0]\,\mathrm dt$。净电荷 $\int I\,\mathrm dt$ 与这两个方向积分不是同一量，在电流变号时尤其不能互换。

如果提供的是电流密度，则必须有可追溯的电极面积才能转换为 A。本版标准数据合同只接受已明确为 A 的电化学电流，不从文件名或软件默认面积推断真实电极面积。

\section{原始电化学时间轴上的积分}
对于相邻点 $(t_i,I_i),(t_{i+1},I_{i+1})$，假设区间内线性插值。如果 $I_iI_{i+1}<0$，零点位于
\begin{equation}
 t_*=t_i+(t_{i+1}-t_i)\frac{-I_i}{I_{i+1}-I_i}.
\end{equation}
先在零点分段，再分别计算两个三角形面积，可以精确积分该分段线性函数的正、负部分。若先把端点截为正值或负值再对整个区间用梯形法，会在变号处引入偏差。

积分窗口边界落在采样点之间时也按同一线性约定插入边界值。程序不把原始 EC 电流改采样到 MS 网格，以免损失短时电流信息。数据缺口、真实充电电流和未观测电流过程仍可能使积分解释产生不确定性，不能由数值积分自动修复。

\section{标定记录由操作者提供}
\code{ms_quantify} 要求工作区 \code{calibrations/} 中存在由操作者提供的记录。工具目录没有“生成实验标定”“批准标定”或“将 verified 改为真”的操作。记录必须包含物种标签、仪器与实验条件标识、有效灵敏度及单位、电子数、参考通道、响应矩阵和核的 SHA-256、审核者，以及是否为合成标定。

这些字段可防止误用，但不是数字签名。持有同一文件系统写权限的人仍能修改记录；部署方若需要法规级审计，应另行采用签名、身份权限、只追加日志和访问控制。本地哈希只用于检查输入和结果是否一致。

\section{定量前的逐项检查}
绝对定量路径要求源项求解数值收敛、当前谱响应满列秩且未触发近共线警告、使用带传输核的模型、质谱信号单位为 A，并且有原始电化学电流。保守的谱诊断门槛可能拒绝某些有额外时间信息可分离的案例；本版选择拒绝而不是自动论证此类可辨识性。

记录中的仪器、条件、合成状态、物种标签、响应及核指纹必须与当前任务匹配；参考通道归一化、单位核质量、正灵敏度和有效积分窗也必须满足要求。完整响应窗由操作者明确确认。若源项为负，程序不会先裁零再给出看似合理的产量。

即使全部门槛通过，结果仍写为“给定响应、传输、标定和电荷定义下的条件定量”。$\mathrm{FE}>1$ 会触发警告，而不是被截为 100\%。不满足门槛时返回可定位的错误，不输出一个默认 FE。

\section{误差预算}
在忽略相关性的一阶近似下，若 $N=A_x/\kappa$，则
\begin{equation}
\left(\frac{u(N)}N\right)^2\simeq
\left(\frac{u(A_x)}{A_x}\right)^2+
\left(\frac{u(\kappa)}\kappa\right)^2.
\end{equation}
实际 $A_x$ 还受背景、传输核、时间窗口和谱矩阵误差影响，这些误差可能相关。固定模型的 Bootstrap 只能覆盖已假定的噪声来源，不能替代完整误差预算。本版 MCP 定量结果不自动声称给出了包含所有系统误差的置信区间。

\section{分电流与 Tafel 分析的边界}
经验证的产物速率可转换为 $i_j=z_jFr_j$；再以明确面积得到分电流密度。Tafel 关系考察过电位与分电流密度对数之间的关系，例如
\begin{equation}
\eta=a+b\log_{10}|j_j|.
\end{equation}
它不是质谱信号对时间的斜率，也不能仅因电位在扫描就从延迟的 MS 峰形直接读出动力学斜率。电位参考、iR 处理、稳态和传质区间均需独立判断。本版 MCP 没有“自动确认 Tafel 区间”的工具，报告亦不把该功能列为已完成。
